\documentclass[
    a4paper,
    11pt,
    bibliography=totoc,
    listof=totoc,
    parskip=half,
    headings=normal
]{scrbook}

% ------------------------------------------------------------------------------
% Packages
% ------------------------------------------------------------------------------
\usepackage[utf8]{inputenc}
\usepackage[T1]{fontenc}
\usepackage[ngerman]{babel}     % German language support

\usepackage[autostyle=true]{csquotes} % Context sensitive quotes
\usepackage{lmodern}            % Modern font
\usepackage{graphicx}           % Graphics support
\usepackage{caption}            % Captions
\usepackage{xcolor}             % Color support
\usepackage{scrlayer-scrpage}   % Header and footer control

% Bibliography setup
\usepackage[
    backend=biber,
    style=verbose-ibid,
    sorting=nty
]{biblatex}
\addbibresource{main.bib}

% Hyperlinks (should be loaded last)
\usepackage[
    hidelinks,
    breaklinks=true,
    pdfauthor={Jan Kaczmarek},
    pdftitle={Agiles Projektmanagement},
    pdfsubject={Kommunikationsdesign}
]{hyperref}

% ------------------------------------------------------------------------------
% Document Start
% ------------------------------------------------------------------------------
\begin{document}

\frontmatter

% Title Page
\begin{titlepage}
  \centering
  \vspace*{2cm}
  {\scshape\LARGE Agiles Projektmanagement\par}
  \vspace{0.5cm}
  {\large Einsendeaufgabe B-APMB01-XX3-N01\par}
  \vspace{2cm}
  {\Large\itshape Jan Kaczmarek\par}
  \vfill
  Studiengang Kommunikationsdesign\par
  \today\par
\end{titlepage}

\tableofcontents

\mainmatter

% ==============================================================================
% 1. Projektmanagement (40 Pkt.)
% ==============================================================================
\chapter{Projektmanagement}

% ------------------------------------------------------------------------------
% 1a) Projektbeschreibung, Expertise, Phasenplan (15 Pkt.)
% ------------------------------------------------------------------------------
\section{Projektbeschreibung, Expertise und Phasenplan}

\subsection{Projektbeschreibung}
Das vorliegende Vorhaben, der Einführung eines Softwaretools zur Digitalisierung von
Geschäftsprozessen in einer mittelständischen Tischlerei, erfüllt alle wesentlichen
Kriterien nach DIN~69\,901.
\begin{itemize}
  \item Einmaligkeit
  \item Zielvorgabe
  \item Begrenzung der verfügbaren Ressourcen
  \item Abgrenzung gegenüber anderen Vorhaben
  \item Projektspezifische Organisation
\end{itemize}

\subsubsection{Einmaligkeit}
Die Digitalisierung der Geschäftsprozesse ist für die Tischlerei ein einmaliges,
neuartiges Vorhaben, das sich grundlegend von der täglichen Routinearbeit im
Handwerksbetrieb unterscheidet. Wie im Studienheft PRMA-H erläutert, hat ein Projekt
im Gegensatz zur immer in gleicher Weise ablaufenden Routinearbeit einen hohen
Neuigkeitsgrad.

\subsubsection{Zielvorgabe}
Die Anpassung und Einführung einer am Markt verfügbaren Software, um die bisher
analog ablaufenden Verwaltungsprozesse der Tischlerei zu digitalisieren. Das
Produkt -- im Sinne des Projektmanagements der Output des Projektes -- ist die
erfolgreich implementierte und an die betrieblichen Bedürfnisse angepasste Software.
Als sekundäre Ziele ergeben sich (a) die Zufriedenstellung der Kunden, welche sich
zuvor über mangelnde Digitalisierung beschwert haben, und (b) die Effizienzsteigerung
der betrieblichen Abläufe. Daraus~(b) ergeben sich potenziell positive Nebeneffekte,
wie ein gesteigertes Betriebsklima.

\subsubsection{Begrenzung der verfügbaren Mittel}

\paragraph{Finanzielle Begrenzung}
Kaufpreis der Software zuzüglich 10\,\% für die Anpassung sowie weitere 10\,\% für
allgemeine Overhead-Kosten.

\paragraph{Personelle Begrenzung}
Personell ist das Projektteam begrenzt, da die Mitarbeitenden neben ihren
Kernaufgaben im Tischlereibetrieb mitwirken müssen.

\paragraph{Zeitliche Begrenzung}
Eine zeitliche Begrenzung ergibt sich aus der finanziellen sowie der personellen
Begrenzung. Eine rasche Umsetzung ist geboten, damit die Overhead-Kosten nicht ihr
Limit übersteigen.

\subsubsection{Abgrenzung und projektspezifische Organisation}
Die Einführung der Software grenzt sich klar von den üblichen Abläufen in der
Tischlerei ab. Neben einer inhaltlich anderen Domäne müssen die regulären Prozesse
im Betrieb parallel weiterlaufen. Diese im Betrieb übliche Linienorganisation ist
für die Steuerung eines solchen Vorhabens nicht geeignet. Das Projekt lässt sich als
Organisationsprojekt mit Investitionscharakter einordnen, da sowohl organisatorische
Veränderungen (Prozessdigitalisierung) als auch die Beschaffung und Anpassung der
Software im Vordergrund stehen.

\subsection{Notwendige Expertise der Mitarbeitenden}
Das Fallbeispiel liefert keine detaillierten Informationen über die genaue Anzahl der
Mitarbeitenden. Daher wird im Folgenden davon ausgegangen, dass die Tischlerei neben
den Beschäftigten in der Werkstatt auch über Verwaltungspersonal verfügt. Zudem wird
angenommen, dass sich das Kollegium aus verschiedenen Altersgruppen zusammensetzt,
die sich hinsichtlich ihrer Vorkenntnisse im Umgang mit digitalen Produkten sowie
ihrer Adaptionsfreudigkeit unterscheiden.

Eine in der \textbf{Verwaltung angestellte Person} kennt die bestehenden Prozesse
(Angebots-, Rechnungsstellung, Bestätigung, Kommunikation~\dots) im Detail und kann
bewerten, wie diese von der Software abgebildet werden müssen.

Eine in der \textbf{Werkstatt angestellte Person} vertritt die Perspektive der
handwerklich tätigen Mitarbeitenden, die zukünftig ebenfalls mit dem System arbeiten
werden (z.\,B. Materialbestellungen, Auftragsdaten abrufen). Es ist ratsam, hier die
Expertise \textbf{einer jüngeren und einer älteren Person} zu unterscheiden. Während
eine Auszubildende oder ein junger Geselle technisch affiner und potenziell weniger
voreingenommen ist, bringt eine bereits lange in der Tischlerei arbeitende Person
tiefes Verständnis für das Handwerk sowie die \enquote{gelebten} Prozesse mit.
Außerdem können Akzeptanzprobleme durch die Mitbestimmung und -gestaltung bereits
während der Integration des IT-Systems ausgeräumt werden.

Der \textbf{externe Kundenbetreuer} der Software-Firma bringt Fachkenntnis zur
Software, zu Anpassungsmöglichkeiten und zur Implementierung mit. Er fungiert als
Schnittstelle zum Softwareanbieter, verfügt jedoch vermutlich nur über wenig Wissen
über das Handwerk, was zu Verständnisproblemen führen kann.

Der \textbf{externe Software-Entwickler} der Software-Firma ist für die technische
Anpassung des Tools an die spezifischen Anforderungen der Tischlerei zuständig. Er
verfügt ebenfalls vermutlich nur über wenig Wissen über das Handwerk.

Der \textbf{Projektleiter (Meister)} bringt umfassendes Prozesswissen mit und kennt
die Schwachstellen, die ihn dazu veranlasst haben, das Projekt anzustoßen. Gemeinsam
mit seiner Digitalaffinität qualifiziert ihn dies für die Leitung des Projekts. Seine
Aufgabe ist es, die Mitarbeitenden sach- und projektkundig anzuleiten und die
Projektaufgabe zu strukturieren. Darüber hinaus dient er als Schnittstelle zur
Softwarefirma.

Optional kann der \textbf{Chef} eingebunden werden. Es wird angenommen, dass er
sowohl über handwerkliche Expertise als auch über mindestens grundlegende
betriebswirtschaftliche Kenntnisse verfügt. Damit wäre er ideal als Mitglied in
einem Lenkungsausschuss geeignet, um die Durchführung zu überwachen, Prioritäten zu
setzen und die Einhaltung des Budgets sicherzustellen.

\subsection{Phasenplan mit Meilensteinen}


% ------------------------------------------------------------------------------
% 1b) Weitere Instrumente zur Projektplanung (10 Pkt.)
% ------------------------------------------------------------------------------
\section{Weitere Instrumente zur Projektplanung}

\subsection{Überblick über Planungsinstrumente}

\subsection{Geeignete Instrumente für dieses Projekt}


% ------------------------------------------------------------------------------
% 1c) Fortschrittsmonitoring (15 Pkt.)
% ------------------------------------------------------------------------------
\section{Fortschrittsmonitoring}

\subsection{Methoden zur Beurteilung des Projektfortschritts}

\subsection{Geeignete Methoden unter Berücksichtigung des Budget-Fokus}


% ==============================================================================
% 2. Psychologie des Projektmanagements (20 Pkt.)
% ==============================================================================
\chapter{Psychologie des Projektmanagements}

% ------------------------------------------------------------------------------
% 2a) Gründe für Projektorganisation (10 Pkt.)
% ------------------------------------------------------------------------------
\section{Gründe für die Projektorganisation}

\subsection{Allgemeine Gründe}

\subsection{Relevanz für das Fallbeispiel}


% ------------------------------------------------------------------------------
% 2b) Visualisierung in Projektbesprechungen (10 Pkt.)
% ------------------------------------------------------------------------------
\section{Gestaltung von Projektbesprechungen und Visualisierung}

\subsection{Vor- und Nachteile der Visualisierung}

\subsection{Geeignete Medien und Auswahl für dieses Projekt}


% ==============================================================================
% 3. Agiles Projektmanagement (40 Pkt.)
% ==============================================================================
\chapter{Agiles Projektmanagement}

% ------------------------------------------------------------------------------
% 3a) Klassisch vs. Agil (15 Pkt.)
% ------------------------------------------------------------------------------
\section{Klassische vs.\ agile Projektmanagementansätze}

\subsection{Vor- und Nachteile des klassischen Ansatzes}

\subsection{Vor- und Nachteile des agilen Ansatzes}

\subsection{Entscheidung für dieses Projekt}


% ------------------------------------------------------------------------------
% 3b) Projektplanung nach Scrum (15 Pkt.)
% ------------------------------------------------------------------------------
\section{Projektplanung nach Scrum}

\subsection{Bausteine von Scrum}

\subsection{Projektablauf}


% ------------------------------------------------------------------------------
% 3c) Verteiltes Team und Kommunikation (10 Pkt.)
% ------------------------------------------------------------------------------
\section{Kommunikation im verteilten Projektteam}

\subsection{Technische Mittel}

\subsection{Besonders wichtige Aspekte der Zusammenarbeit}


\backmatter

% ------------------------------------------------------------------------------
% Bibliography
% ------------------------------------------------------------------------------
\printbibliography

\end{document}